\makeatletter
\@ifundefined{isChecklistMainFile}{
  % We are compiling a standalone document
  \newif\ifreproStandalone
  \reproStandalonetrue
}{
  % We are being \input into the main paper
  \newif\ifreproStandalone
  \reproStandalonefalse
}
\makeatother

\ifreproStandalone
\documentclass[letterpaper]{article}
\usepackage[submission]{aaai2027}
\setlength{\pdfpagewidth}{8.5in}
\setlength{\pdfpageheight}{11in}
\frenchspacing

\begin{document}
\fi
\setlength{\leftmargini}{20pt}
\makeatletter\def\@listi{\leftmargin\leftmargini \topsep .5em \parsep .5em \itemsep .5em}
\def\@listii{\leftmargin\leftmarginii \labelwidth\leftmarginii \advance\labelwidth-\labelsep \topsep .4em \parsep .4em \itemsep .4em}
\def\@listiii{\leftmargin\leftmarginiii \labelwidth\leftmarginiii \advance\labelwidth-\labelsep \topsep .4em \parsep .4em \itemsep .4em}\makeatother

\setcounter{secnumdepth}{0}
\renewcommand\thesubsection{\arabic{subsection}}
\renewcommand\labelenumi{\thesubsection.\arabic{enumi}}

\newcounter{checksubsection}
\newcounter{checkitem}[checksubsection]

\newcommand{\checksubsection}[1]{%
  \refstepcounter{checksubsection}%
  \paragraph{\arabic{checksubsection}. #1}%
  \setcounter{checkitem}{0}%
}

\newcommand{\checkitem}{%
  \refstepcounter{checkitem}%
  \item[\arabic{checksubsection}.\arabic{checkitem}.]%
}
\newcommand{\question}[2]{\normalcolor\checkitem #1 #2 \color{blue}}
\newcommand{\ifyespoints}[1]{\makebox[0pt][l]{\hspace{-15pt}\normalcolor #1}}

\section*{Reproducibility Checklist}

\vspace{1em}
\hrule
\vspace{1em}

% The questions start here

\checksubsection{General Paper Structure}
\begin{itemize}

\question{Includes a conceptual outline and/or pseudocode description of AI methods introduced}{(yes/partial/no/NA)}
NA -- this paper introduces no new AI method; Section 3 (Method) gives a full equation-level description of the existing Byz-Clip21-SGD2M algorithm under test, for clarity.

\question{Clearly delineates statements that are opinions, hypothesis, and speculation from objective facts and results}{(yes/no)}
yes

\question{Provides well-marked pedagogical references for less-familiar readers to gain background necessary to replicate the paper}{(yes/no)}
no

\end{itemize}
\checksubsection{Theoretical Contributions}
\begin{itemize}

\question{Does this paper make theoretical contributions?}{(yes/no)}
no

	\ifyespoints{\vspace{1.2em}If yes, please address the following points:}
        \begin{itemize}

	\question{All assumptions and restrictions are stated clearly and formally}{(yes/partial/no)}
	NA

	\question{All novel claims are stated formally (e.g., in theorem statements)}{(yes/partial/no)}
	NA

	\question{Proofs of all novel claims are included}{(yes/partial/no)}
	NA

	\question{Proof sketches or intuitions are given for complex and/or novel results}{(yes/partial/no)}
	NA

	\question{Appropriate citations to theoretical tools used are given}{(yes/partial/no)}
	yes

	\question{All theoretical claims are demonstrated empirically to hold}{(yes/partial/no/NA)}
	NA

	\question{All experimental code used to eliminate or disprove claims is included}{(yes/no/NA)}
	NA

	\end{itemize}
\end{itemize}

\checksubsection{Dataset Usage}
\begin{itemize}

\question{Does this paper rely on one or more datasets?}{(yes/no)}
yes

\ifyespoints{If yes, please address the following points:}
\begin{itemize}

	\question{A motivation is given for why the experiments are conducted on the selected datasets}{(yes/partial/no/NA)}
	yes -- MNIST replicates the source paper's own validation dataset; CIFAR-10 is added specifically to test H1 (whether the source paper's sub-Gaussian-noise premise and empirical robustness generalize past MNIST scale).

	\question{All novel datasets introduced in this paper are included in a data appendix}{(yes/partial/no/NA)}
	NA -- no novel datasets are introduced.

	\question{All novel datasets introduced in this paper will be made publicly available upon publication of the paper with a license that allows free usage for research purposes}{(yes/partial/no/NA)}
	NA

	\question{All datasets drawn from the existing literature (potentially including authors' own previously published work) are accompanied by appropriate citations}{(yes/no/NA)}
	yes

	\question{All datasets drawn from the existing literature (potentially including authors' own previously published work) are publicly available}{(yes/partial/no/NA)}
	yes

	\question{All datasets that are not publicly available are described in detail, with explanation why publicly available alternatives are not scientifically satisficing}{(yes/partial/no/NA)}
	NA

\end{itemize}
\end{itemize}

\checksubsection{Computational Experiments}
\begin{itemize}

\question{Does this paper include computational experiments?}{(yes/no)}
yes

\ifyespoints{If yes, please address the following points:}
\begin{itemize}

	\question{This paper states the number and range of values tried per (hyper-) parameter during development of the paper, along with the criterion used for selecting the final parameter setting}{(yes/partial/no/NA)}
	yes -- every hyperparameter-probe stage reports every value tried, including configurations that did not work, and states the selection criterion (highest clean-training accuracy within the fixed round budget).

	\question{Any code required for pre-processing data is included in the appendix}{(yes/partial/no)}
	no -- released via the accompanying code repository rather than a paper appendix (see code/data availability statement).

	\question{All source code required for conducting and analyzing the experiments is included in a code appendix}{(yes/partial/no)}
	no -- released via the accompanying code repository rather than a paper appendix.

	\question{All source code required for conducting and analyzing the experiments will be made publicly available upon publication of the paper with a license that allows free usage for research purposes}{(yes/partial/no)}
	[TODO: confirm with authors before submission -- not yet decided in this pilot.]

	\question{All source code implementing new methods have comments detailing the implementation, with references to the paper where each step comes from}{(yes/partial/no)}
	partial -- the core algorithm implementation (src/byz\_clip21\_sgd2m.py) is commented line-by-line against the source paper's pseudocode; auxiliary scripts are less exhaustively annotated.

	\question{If an algorithm depends on randomness, then the method used for setting seeds is described in a way sufficient to allow replication of results}{(yes/partial/no/NA)}
	yes -- every run fixes \texttt{torch.manual\_seed} and \texttt{np.random.seed} to an explicit integer seed recorded in that run's output JSON and in the paper's Reproducibility section seed list.

	\question{This paper specifies the computing infrastructure used for running experiments (hardware and software), including GPU/CPU models; amount of memory; operating system; names and versions of relevant software libraries and frameworks}{(yes/partial/no)}
	yes -- CPU-only (no GPU/CUDA available), 8 logical cores, Windows, Python 3.11.9, PyTorch 2.13.0+cpu, torchvision 0.28.0+cpu, NumPy 2.4.6, SciPy 1.17.1.

	\question{This paper formally describes evaluation metrics used and explains the motivation for choosing these metrics}{(yes/partial/no)}
	yes -- final test accuracy for training runs; Hill tail-index alpha, excess kurtosis, KS statistic vs.\ N(0,1), and Gaussian-/Student-t-referenced exceedance ratios for the sub-Gaussianity measurements, each motivated as a distinct, complementary probe of tail heaviness.

	\question{This paper states the number of algorithm runs used to compute each reported result}{(yes/no)}
	yes -- this is a first-class design requirement of this project; every table states its exact seed count, and under-powered tables are explicitly flagged rather than silently presented as if fully powered.

	\question{Analysis of experiments goes beyond single-dimensional summaries of performance (e.g., average; median) to include measures of variation, confidence, or other distributional information}{(yes/no)}
	yes -- MNIST tables (n=10 or n=5), the CIFAR-10 tau-search confirm phase (n=5), and CIFAR-10's C1/C2/C3 main-sweep tables (n=10/n=5/n=10, brought up from the original pilot's n=2--3 in a follow-up compute pass) are all reported with std, and min/max where tabulated (see Limitations for the one remaining n=1 case, the CIFAR-10 Stage-A hp-probe).

	\question{The significance of any improvement or decrease in performance is judged using appropriate statistical tests (e.g., Wilcoxon signed-rank)}{(yes/partial/no)}
	partial -- paired Wilcoxon signed-rank tests (scipy.stats.wilcoxon, matched by seed) are now reported for the paper's main comparative claims, in a dedicated significance-testing table in the Results section; no correction for multiple comparisons is applied across the 11 tests run, which is a disclosed limitation (see Limitations), not an oversight.

	\question{This paper lists all final (hyper-)parameters used for each model/algorithm in the paper's experiments}{(yes/partial/no/NA)}
	yes -- every run's full configuration (including seed, $\gamma$, $\tau$, $\beta$, $\hat\beta$, $\varepsilon$, RAgg choice, ablation flag) is saved verbatim in that run's output JSON under \texttt{results/}, and every table in the paper traces to specific files.

\end{itemize}
\end{itemize}
\ifreproStandalone
\end{document}
\fi
